\documentclass[11pt,a4paper]{article}

\usepackage[margin=2.2cm]{geometry}
\usepackage{booktabs}
\usepackage{multirow}
\usepackage{array}
\usepackage{amsmath,amssymb}
\usepackage{xcolor}
\usepackage{hyperref}
\usepackage{parskip}
\usepackage{lmodern}
\usepackage{fancyvrb}

\hypersetup{colorlinks=true, linkcolor=black, urlcolor=blue}

\title{Phase 9 — Composability Test\\[2pt]
       Entropy regularisation vs. Batch-Norm, Dropout, Weight-Decay}
\author{thesis-IV working notes (HyMeKo)}
\date{2026-04-26}

\newcommand{\code}[1]{\texttt{\detokenize{#1}}}

\begin{document}
\maketitle

\section{Purpose}

Phases 6--8 demonstrated that
\code{scalar_entropy_normalized} (Path~B) and \code{entropy_target}
($H^*=0.5$, Path~A) universally Pareto-dominate the original
\code{scalar_entropy} regulariser across nine datasets and five
architectures. \textbf{Open question for the manuscript:} is the new
regulariser \emph{competitive} with mainstream techniques alone, and
\emph{complementary} when stacked on top of them?

Phase 9 answers that by running paired comparisons
\code{baseline} vs \code{scalar_entropy_normalized} under each of five
model + optimiser settings:

\begin{center}
\small
\begin{tabular}{@{}llp{8cm}@{}}
\toprule
ID & Variant         & Modification \\
\midrule
V1 & \code{vanilla}      & base model, Adam(lr=$10^{-3}$). \\
V2 & \code{dropout}      & + \code{nn.Dropout}($p{=}0.5$) after each ReLU \\
V3 & \code{batch_norm}   & + \code{nn.BatchNorm1d} after each Linear \\
V4 & \code{weight_decay} & vanilla model, Adam($\text{wd}{=}10^{-4}$) \\
V5 & \code{full_stack}   & BN + Dropout + Adam($\text{wd}{=}10^{-4}$) \\
\bottomrule
\end{tabular}
\end{center}

The pair \code{baseline} vs \code{scalar_entropy_normalized} within each
variant tells us whether entropy adds value \emph{on top of} the existing
regularisation. Cross-variant comparisons of mean accuracies tell us
whether entropy alone is competitive.

\section{Configuration matrix}

Five variants $\times$ four datasets = \textbf{20 paired comparisons}
(40 individual training runs). Estimated wall-clock $\sim$11~h.
Auto-fires after \texttt{Views PH8 suite finished} marker.

\begin{center}
\small
\begin{tabular}{@{}llrr@{}}
\toprule
Dataset/arch & Original sign & Seeds & Epochs \\
\midrule
\code{spirals}        & strongly + & 100 & 50 \\
\code{circles}        & sig.\ $-$  & 100 & 50 \\
\code{mnist_small}    & + (anchor) & 33  & 15 \\
\code{mnist_capsnet}  & sig.\ $-$  & 33  & 10 \\
\bottomrule
\end{tabular}
\end{center}

\section{Per-cell bash commands}

For each (dataset, variant, seeds, epochs) configuration the harness
is invoked exactly as below. The base shell script is
\code{python/benches/thesis_iv_hard/run_overnight_views_ph9.sh}.

\subsection{Block A: synthetic 2-D (spirals + circles)}

\begin{Verbatim}[fontsize=\small,frame=single]
# V1 vanilla
python3 .../run_benchmark.py --datasets spirals \
    --arms baseline scalar_entropy_normalized \
    --seeds 100 --epochs 50 --lam 0.1 --reg-every-n 10 --view dataflow

# V2 dropout
python3 .../run_benchmark.py --datasets spirals_drop \
    --arms baseline scalar_entropy_normalized \
    --seeds 100 --epochs 50 --lam 0.1 --reg-every-n 10 --view dataflow

# V3 batch norm
python3 .../run_benchmark.py --datasets spirals_bn \
    --arms baseline scalar_entropy_normalized \
    --seeds 100 --epochs 50 --lam 0.1 --reg-every-n 10 --view dataflow

# V4 weight decay
python3 .../run_benchmark.py --datasets spirals \
    --arms baseline scalar_entropy_normalized \
    --seeds 100 --epochs 50 --lam 0.1 --reg-every-n 10 --view dataflow \
    --weight-decay 1e-4

# V5 full stack
python3 .../run_benchmark.py --datasets spirals_full \
    --arms baseline scalar_entropy_normalized \
    --seeds 100 --epochs 50 --lam 0.1 --reg-every-n 10 --view dataflow \
    --weight-decay 1e-4
\end{Verbatim}

The same five commands repeat with \code{--datasets circles*} for the
second synthetic block.

\subsection{Block B: MNIST plain MLP}

\begin{Verbatim}[fontsize=\small,frame=single]
# V1, V2, V3
python3 .../run_benchmark.py --datasets {mnist_small | _drop | _bn} \
    --arms baseline scalar_entropy_normalized \
    --seeds 33 --epochs 15 --lam 0.1 --reg-every-n 10 --view dataflow

# V4 (vanilla + WD)
python3 .../run_benchmark.py --datasets mnist_small \
    --arms baseline scalar_entropy_normalized \
    --seeds 33 --epochs 15 --lam 0.1 --reg-every-n 10 --view dataflow \
    --weight-decay 1e-4

# V5 full stack
python3 .../run_benchmark.py --datasets mnist_small_full \
    --arms baseline scalar_entropy_normalized \
    --seeds 33 --epochs 15 --lam 0.1 --reg-every-n 10 --view dataflow \
    --weight-decay 1e-4
\end{Verbatim}

\subsection{Block C: CapsMLP MNIST}

\begin{Verbatim}[fontsize=\small,frame=single]
# V1, V2, V3
python3 .../run_benchmark.py --datasets {mnist_capsnet | _drop | _bn} \
    --arms baseline scalar_entropy_normalized \
    --seeds 33 --epochs 10 --lam 0.1 --reg-every-n 10 --view dataflow

# V4 (vanilla CapsMLP + WD)
python3 .../run_benchmark.py --datasets mnist_capsnet \
    --arms baseline scalar_entropy_normalized \
    --seeds 33 --epochs 10 --lam 0.1 --reg-every-n 10 --view dataflow \
    --weight-decay 1e-4

# V5 full stack
python3 .../run_benchmark.py --datasets mnist_capsnet_full \
    --arms baseline scalar_entropy_normalized \
    --seeds 33 --epochs 10 --lam 0.1 --reg-every-n 10 --view dataflow \
    --weight-decay 1e-4
\end{Verbatim}

\section{Hypotheses for the paper claim}

\begin{enumerate}
  \item \textbf{Composability hypothesis (the killer finding).}
        The paired $\Delta$ from adding \code{scalar_entropy_normalized}
        on top of V2--V5 is \emph{at least as positive} as on V1
        (vanilla). If yes, the regulariser \textbf{stacks} with existing
        techniques without conflict, which is the core practical claim.
  \item \textbf{Standalone competitiveness.}
        The absolute mean accuracy of V1+entropy is comparable to or
        better than the absolute mean accuracy of V2--V5 standalone
        (i.e.\ baseline arm under each variant). If yes,
        \code{scalar_entropy_normalized} on its own is a credible
        single-tool regulariser.
  \item \textbf{Failure-mode mitigation.}
        On the negative datasets (\code{circles}, \code{mnist_capsnet}),
        none of dropout / BN / WD alone neutralises the failure mode of
        the original \code{scalar_entropy}. If only V2--V5 \emph{with}
        entropy do, the regulariser is a unique tool, not a
        substitutable one.
  \item \textbf{Variance reduction.}
        The variance-reduction signature should survive — and ideally
        compound — when entropy is added on top of dropout / BN / WD.
\end{enumerate}

\section{Acceptance criteria}

For the paper to claim ``composability with mainstream regularisers'',
at least 3 of 4 dataset $\times$ variant cells in V5 (full stack) must
show $\Delta \geq 0$ from adding \code{scalar_entropy_normalized}, and
at least 2 of those must reach significance ($p < 0.05$).

For the paper to claim ``failure-mode mitigation'', V1, V2, V3, V4
under \code{baseline} on \code{circles} and \code{mnist_capsnet} must
all retain a directional negative effect; while V1--V5 under
\code{scalar_entropy_normalized} must all be null or positive.

\section{After phase 9}

Refresh \code{RESULTS_VIEWS_SUITE.md} via the existing aggregator. The
new rows will appear under their dataset name. Phase 9 cells are
identifiable by the \code{*_drop}, \code{*_bn}, \code{*_full} dataset
suffixes and (where applicable) by the non-zero \texttt{weight\_decay}
in the per-run config printout.
\end{document}
